\item Un átomo aislado de cierto elemento emite luz de $520\text{ nm}$ de longitud de onda cuando el átomo cae de su quinto estado excitado a su segundo estado excitado. El átomo emite un fotón de $410\text{ nm}$ de longitud de onda cuando cae de su sexto estado excitado a su segundo estado excitado. Encuentre la longitud de onda de la luz radiada cuando el átomo hace una transición de su sexto a su quinto estado excitado.